% Complete Section 8.5, original PDF 111-112 / printed 101-102.
% block: parabolic-hyperbolic-heading page:111 type:section
\section{{抛物型与双曲型问题}}\label{sec:parabolic-hyperbolic}

% block: hodge-heat-equation page:111 type:ordinary-paragraph
{除 Hodge Laplace 方程 $(d\delta+\delta d)u=f$ 外, FEEC 还可用于构造并求解相关的抛物型方程, 即求满足 \emph{Hodge 热方程}的\emph{含时}微分形式 $u:[0,T]\to\Lambda^k$}
\begin{equation}\label{eq:hodge-heat-equation}
u_t+(\delta d+d\delta u)=f,
\end{equation}
{并附有\emph{初始条件} $u(0)=u_0$, 以及例如自然边界条件 $\operatorname{tr}\star u=0$ 和 $\operatorname{tr}\star du=0$. 与椭圆情形一样, 弱形式引入辅助变量 $\sigma=\delta u$, 并由下式确定 $(\sigma,u):[0,T]\to H\Lambda^{k-1}\times H\Lambda^k$}
\begin{equation}\label{eq:hodge-heat-weak-formulation}
\begin{aligned}
\langle\sigma,\tau\rangle-\langle d\tau,u\rangle&=0,&\tau&\in H\Lambda^{k-1},\ t\in[0,T],\\
\langle u_t,v\rangle+\langle d\sigma,v\rangle+\langle du,dv\rangle&=\langle f,v\rangle,&v&\in H\Lambda^k,\ t\in[0,T],
\end{aligned}
\end{equation}
{并附上初始条件. 注意, \eqref{eq:hodge-heat-weak-formulation} 中的方程必须在每个时刻 $t$ 成立. 与 \eqref{eq:general-hodge-weak-formulation} 相比, \eqref{eq:hodge-heat-weak-formulation} 多了一个含时间导数 $u_t$ 的项, 但与椭圆情形不同, 即使存在调和形式, 也不需要在抛物型问题的弱形式中显式计入它们.}

% block: hodge-heat-analysis pages:111,112 type:ordinary-paragraph
{[\BookNumber{10}, Theorem 4.4] 在 FEEC 框架下讨论了 Hodge 热方程. 其中说明, 如何将已得到的 Hodge Laplacian 结果与 Hille–Yosida–Phillips 理论结合, 建立抛物型问题的适定性. 该文继而分析 Galerkin 方法, 包括半离散方法, 以及采用简单时间步进格式的全离散方法. 基于适当的椭圆投影得到了最优阶误差估计, 该投影的性质同样来自椭圆情形.}

% block: hodge-wave-equation page:112 type:ordinary-paragraph
{还可考虑 \emph{Hodge 波动方程}, 它是二阶双曲型方程, 通过下述微分方程确定含时 $k$-形式 $u:[0,T]\to\Lambda^k$}
\begin{equation}\label{eq:hodge-wave-equation}
u_{tt}+(\delta d+d\delta)u=f.
\end{equation}

% block: hodge-wave-first-order-system page:112 type:ordinary-paragraph
{当然, 现在需要两个初始条件, $u(0)=u_0$ 和 $u_t(0)=v_0$. 为求解此问题, 引入三个辅助变量: $\sigma=\delta u$, $v=u_t$ 和 $\beta=du$, 它们分别是次数为 $k-1$, $k$ 和 $k+1$ 的含时微分形式. 于是可导出强形式方程}
\[
\begin{aligned}
\sigma_t&=\delta u_t=\delta v,\\
v_t&=u_{tt}=-(\delta d+d\delta)u+f=-\delta\beta-d\sigma+f,\\
\beta_t&=du_t=dv,
\end{aligned}
\]
{或写成矩阵形式,}
\[
\frac d{dt}\begin{pmatrix}\sigma\\v\\\beta\end{pmatrix}=
\begin{pmatrix}0&\delta&0\\-d&0&-\delta\\0&d&0\end{pmatrix}
\begin{pmatrix}\sigma\\v\\\beta\end{pmatrix}+\begin{pmatrix}0\\f\\0\end{pmatrix},
\]
{这揭示了对称双曲系统的结构 (关键在于右端的矩阵算子是反对称的). 初始数据取为 $\sigma(0)=\delta u_0$, $v(0)=v_0$ 和 $\beta(0)=du_0$. 注意, 原变量 $u$ 不是此系统的一部分, 但可通过积分 $u_t=v$ 得到, 即}
\[
u(t)=u_0+\int_0^t v(s)\,ds.
\]

% block: hodge-wave-weak-formulation page:112 type:ordinary-paragraph
{现在转向混合弱形式, 求 $(\sigma,v,\beta):[0,T]\to H\Lambda^{k-1}\times H\Lambda^k\times H\Lambda^{k+1}$ 满足}
\begin{equation}\label{eq:hodge-wave-weak-formulation}
\begin{aligned}
\langle\sigma_t,\tau\rangle-\langle v,d\tau\rangle&=0,&\tau&\in H\Lambda^{k-1},\ t\in[0,T],\\
\langle v_t,w\rangle+\langle d\sigma,w\rangle+\langle\beta,dw\rangle&=\langle f,w\rangle,&w&\in H\Lambda^k,\ t\in[0,T],\\
\langle\beta_t,\gamma\rangle-\langle dv,\gamma\rangle&=0,&\gamma&\in H\Lambda^{k+1},\ t\in[0,T].
\end{aligned}
\end{equation}

% block: hodge-wave-wellposedness-energy page:112 type:ordinary-paragraph
{[\BookNumber{63}, Chapter 4] 在 Hilbert 复形框架下建立了此形式的适定性及其 Galerkin 解的收敛性. 与抛物型问题一样, 这些结果利用 Hille–Yosida–Phillips 理论得到; 当不含时间的算子像这里一样是反对称的时, 该理论尤其简单. 同一反对称性也立即给出能量守恒. 也就是说, 在弱形式中取测试函数 $(\tau,w,\gamma)$ 等于时刻 $t$ 的解 $(\sigma,v,\beta)$, 可见}
\[
\|\sigma\|_{L^2}^2+\|v\|_{L^2}^2+\|\beta\|_{L^2}^2
\]
{不随时间改变.}
